\documentclass[11pt,a4paper]{article}
\usepackage{fontspec}
\setmainfont{TeX Gyre Termes}
\usepackage{polyglossia}
\setmainlanguage{english}
\setotherlanguage{portuguese}
\usepackage{amsmath,booktabs,geometry,hyperref}
\geometry{left=3cm,right=3cm,top=2.5cm,bottom=2.5cm}
\hypersetup{colorlinks=true,linkcolor=black,urlcolor=blue}
\title{Listening to structure in electroacoustic performance}
\author{\textbf{Author One} \and \textbf{Author Two}}
\date{Per Musi --- volume and issue to be completed}
\newcommand{\resumo}[1]{\par\smallskip\textbf{Resumo.} #1}
\newcommand{\keywords}[1]{\par\smallskip\textbf{Keywords:} #1}

\begin{document}
\maketitle
\begin{abstract}
This article presents a listening protocol for comparing repeated
electroacoustic performances. It combines annotated listening sessions with a
small timing measure so interpretive observations remain traceable.
\end{abstract}
\resumo{Este artigo apresenta um protocolo de escuta para comparar performances
eletroacústicas repetidas, combinando observações anotadas e uma medida temporal.}
\keywords{electroacoustic music; listening; performance; analysis}

\section{Introduction}
An analytical account should say what was heard, when it was heard, and which
features support the interpretation. This starter keeps those fields visible
while leaving the final musical examples to the author.

\section{Method}
For an annotated event at time $t_i$, we record its distance from the expected
cue $c_i$ as $d_i=|t_i-c_i|$. The mean deviation is
\begin{equation}
  D=\frac{1}{n}\sum_{i=1}^{n}|t_i-c_i|.
\end{equation}
Qualitative notes should accompany the numeric summary rather than being
replaced by it.

\section{Results}
Table~\ref{tab:events} gives an illustrative event log.
\begin{table}[ht]
\centering
\caption{Annotated performance events.}
\label{tab:events}
\begin{tabular}{lrr}
\toprule
Event & Cue (s) & Observed (s)\\
\midrule
Entrance & 42.0 & 42.4\\
Transition & 96.0 & 95.7\\
Release & 151.0 & 151.2\\
\bottomrule
\end{tabular}
\end{table}

\section{Conclusion}
The protocol links a listening claim to a timestamp and an observation. Add
the work's final repertoire, recordings, ethics statement, and editorial
metadata before submission.

\begin{thebibliography}{9}
\bibitem{cook} N. Cook, \emph{Beyond the Score: Music as Performance}, Oxford
University Press, 2013.
\bibitem{clarke} E. Clarke, \emph{Ways of Listening}, Oxford University Press,
2005.
\end{thebibliography}
\end{document}
